\newpage




\subsubsection{$\ln$: Logarithmus Naturalis}

Geben Sie mit Hilfe des natürlichen Logarithmus $\ln()$ an:


$$7^x = 5^{x+2}$$
\TNT{16}{
$$7^x = 5^x \cdot{} 5^2$$
$$\frac{7^x}{5^x} = 5^2$$
$$\left(\frac{7}{5}\right)^x = 5^2$$
$$x = \frac{\ln(5^2)}{\ln\left(\frac{7}{5}\right)}=$$
Wie ist das Resultat anzugeben?
$$x=\frac{2\cdot{}\ln(5)}{\ln(7)-\ln(5)} (\text{exakt}))$$
$$x\approx 9.56654212303371468116 (\text{angenähert})$$
}%% end TNT

\TALS{\olatLinkArbeitsblatt{Exponentialgleichungen [GL\_Ex]}{https://olat.bms-w.ch/auth/RepositoryEntry/6029786/CourseNode/106029166654550}{5.}}

\newpage



\subsubsection{Substitution}

Finden Sie die Lösungsmenge $\lx$:

$$2\cdot{} 5^x - 5^{2x} = 1$$

\TNT{10}{
  Substituiere $u := 5^x$

  $$2u -u^2 = 1$$
  Quadratische Gleichung:
  $$0 = u^2 -2u + 1$$
  $$0 = (u-1)(u-1)$$
  daraus folgt:
  $$u= 1$$
  Rücksubstitution:
  $$1=5^x$$
  Und daraus:
  $$x=0 ; \lx=\{0\}$$

}%% end TNT

\TALS{\olatLinkArbeitsblatt{Exponentialgleichungen [GL\_Ex]}{https://olat.bms-w.ch/auth/RepositoryEntry/6029786/CourseNode/106029166654550}{6.}}

